\footnotesize%
\pstart%
\noindent%
\pagelabel{paradoxonbarometricum_intro}%
Der vorliegende Textkomplex N.\;\ref{paradoxon_intro}, dem editorisch der Titel \textit{Paradoxon de aeris gravitate} zugewiesen wurde, umfasst vier Texte (N.\;\ref{53442} %
bis N.\;\ref{54958}), welche von einem in den Augen der frühneuzeitlichen Naturforschung kontraintuitiven Phänomen handeln:
Das Luftgewicht sinkt bei schelechtem und steigt bei gutem Wetter, wie die Veränderung der Quecksilberhöhe im Barometer zeigt.
In den Aufzeichnungen N.\;\ref{53442} bis N.\;\ref{36805} schlägt Leibniz für dieses Phänomen eine Erklärung vor, die von folgender hydrostatischer Annahme ausgeht:
Ein Körper belaste mit seinem Gewicht eine Flüssigkeit nicht, solange er in ihr versinke.
Das paradoxe Verhalten des Luftgewichts will Leibniz entsprechend begründen:
Bei schlechtem Wetter verringere sich das Luftgewicht, weil die Luftfeuchtigkeit sich in Tropfen verdichte und folglich \glqq absinke\grqq, wodurch die Luftmasse gemäß der besagten Annahme entlastet werde;
nach Verbesserung der Wetterlage erhöhe sich das Luftgewicht wieder, weil die verdichtete Feuchtigkeit erneut verdünste und, mit der Luft vermischt, diese letztere wieder belaste.
Die hydrostatische Annahme, auf der die Erklärung beruht, lässt sich Leibniz zufolge empirisch bestätigen:
Mithilfe einer Waage soll möglich sein, die Entlastung einer flüssigen Masse beim Versinken eines in ihr schwimmenden Körpers zu zeigen.
Ein solches \textit{experimentum mechanicum} (\refpassage{LBr_413_048v_01b}{LBr_413_048v_01b}) -- ein Gedankenexperiment -- stellt Leibniz in den Aufzeichnungen N.\;\ref{53442} bis N.\;\ref{36805} stets in gleicher Weise dar.
Diese drei Texte dürften insgesamt zwischen Mitte September 1696 und dem 19.\;März 1701 abgefasst worden sein (siehe unten).
Später entstanden ist sicherlich die 1715 in den \textit{Ephemerides} der Leopoldina-Akademie veröffentlichten Erläuterung N.\;\ref{54958} \textit{De aetiologia barometri}, in der Leibniz auf die frühere Auseinandersetzung zwischen Bernardino Ramazzini\indexN{\textls{Ramazzini} (Ramazzinus), Bernardino 1633\textendash1714} und Günther Christoph Schelhammer\indexN{\textls{Schelhammer} (Schelhammerus), Günther Christoph 1649\textendash1716} über das paradoxe Verhalten des Luftgewichtes eingeht und seine eigene Erklärung des Phänomens gegen Schelhammers Kritik verteidigt.
Ausschaggebend für die absolute und relative Datierung der einzelnen Texte N.\;\ref{53442} bis N.\;\ref{54958} sind folgende Beobachtungen:
\pend%
\pstart%
(1)
In der titellosen und undatierten Aufzeichnung N.\;\ref{53442}, editorisch \textit{Causa paradoxi cur aer pluvius sit levior} benannt, wird hauptsächlich die oben umgerissene hydrostatische Annahme dargelegt, die aus Leibnizens Sicht die wetterbedingten Veränderungen des Luftgewichts begründen soll (\refpassage{LH_37_04_151_151r-annahme-1}{LH_37_04_151_151r-annahme-2}).
Hierbei greift er stillschweigend einen Gedanken wieder auf, welcher bereits in der von Leibniz eigenhändig auf \textit{März 1677} datierten Aufzeichnung N.\;\ref{60277} im vorliegenden Band geäußert worden war (\refpassage{LH_37_05_018_018r_Gdnkexp-1}{LH_37_05_018_018r_Gdnkexp-2}).
N.\;\ref{53442} ist sicherlich später als N.\;\ref{60277} entstanden.
Dafür spricht zunächst der Umnstand, dass in N.\;\ref{60277} die besagte hydrostatische Annahme noch keine Erklärung des "barometrischen Paradoxons" veranlasst, wie dies hingegen in N.\;\ref{53442} und in beiden späteren Texten N.\;\ref{60267} und \ref{36805} der Fall ist. %
Ferner überliefert der Textträger der Aufzeichnung N.\;\ref{53442} auch eine undatierte Kostenaufstellung, die sich dort bereits vor N.\;\ref{53442} befand und wahrscheinlich von Leibnizens Kutscher und Diener Hans\indexN{Hans, Leibnizens Diener und Kutscher} stammt (vgl. \refpassage{LH_37_04_151_151v_hans-1}{LH_37_04_151_151v_hans-2}). %
Ab wann dieser bei Leibniz angestellt war, ist nicht bekannt.
Erwähnt wird er %
im Leibniz-Briefwechsel insgesamt etwa dreißig Mal zwischem dem 26.\;Juni 1689 (Brief von Johann Daniel Crafft\indexN{\textls{Crafft} (Craftius, Krafft, Kraftius), Johann Daniel 1624\textendash1697} an Leibniz: \textit{LSB} III,\;4 N.\;210, S.\;403.15--18)\cite{02575} und dem 1.\;(11.) März 1698 (Brief von Johann Christoph Balcke\indexN{\textls{Balcke}, Johann Christoph 1661\textendash1730} an Leibniz: \textit{LSB} I,\;15 N.\;30, S.\;43.6)\cite{02576}.
Das umfassende Ausbleiben von Erwähnungen Hansens\indexN{Hans, Leibnizens Diener und Kutscher} im Briefwechsel vor dem Sommer 1689 lässt vermuten, dass er im März 1677 noch kein Dienstverhältnis zu Leibniz hatte.
Aus demselben Grund erweist sich als wenig wahrscheinlich, dass die Aufzeichnung N.\;\ref{53442} vor 1689 entstand.
Vielmehr ist es plausibel, bei N.\;\ref{53442} eine ähnliche Entstehungszeit anzunehmen wie bei der Aufzeichnung N.\;\ref{60267}.
Hierfür spricht hauptsächlich die enge inhaltliche Verwandtschaft beider Texte:
Die in N.\;\ref{53442} knapp, vorläufig und abbrechend formulierte Erklärung wird in N.\;\ref{60267} ausführlich entfaltet und umfassend artikuliert;
der Anfang von N.\;\ref{53442} (\refpassage{LH_37_04_151_151r-annahme-1}{LH_37_04_151_151r-annahme-2}) ist nahezu wörtlich in N.\;\ref{60267} (\refpassage{LH_37_04_149-150_149v-annahme-1}{LH_37_04_149-150_149v-annahme-2}) übernommen;
die jeweiligen Diagramme \bracks{\textit{Fig.\:1}} und \bracks{\textit{Fig.\:2}} weisen eine auffällige Ähnlichkeit auf.
Es ist daher wahrscheinlich, dass N.\;\ref{53442} nicht viel früher als N.\;\ref{60267} verfasst wurde.
Eine Datierung ab 1689 ist jedoch nicht auszuschließen.
\pend%
\pstart%
(2)
Die titellose und undatierte Aufzeichnung N.\;\ref{60267}, editorisch \textit{Paradoxon barometricum ejusque causa} benannt, beginnt mit einem ausdrücklichen und ausführlichen Verweis auf die \textit{Ephemerides barometricae} des italienischen Arztes und Naturforschers Bernardino Ramazzini,\indexN{\textls{Ramazzini} (Ramazzinus), Bernardino 1633\textendash1714} welche zunächst 1695\cite{01516} in Modena\indexO{Modena} und dann 1696\cite{01517} in Augsburg\indexO{Augsburg} (zusammen mit Rudolf Jakob Camerarius'\indexN{\textls{Camerarius}, Rudolf Jakob 1665\textendash1721} \textit{Ephemerides meteorologicae})\cite{01445} veröffentlicht worden waren.
Daraus ist zu schließen, dass N.\;\ref{60267} erst abgefasst wurde, als Leibniz Ramazzinis \textit{Ephemerides barometricae}\cite{01516} gelesen hatte.
Diese Schlussfolgerung trifft auf die Aufzeichnung N.\;\ref{53442} insofern nicht zwingend zu, als dort kein \textit{expliziter} Verweis auf Ramazzinis Abhandlung anzutreffen ist.
In N.\;\ref{53442} (\refpassage{LH_37_04_151_151r_anspiel-1}{LH_37_04_151_151r_anspiel-2}) findet sich allerdings eine Anspielung auf Vorgänger, die Leibnizens Ansatz zwar vorweggenommen, aber unzureichend begründet hätten.
Diese Anspielung dürfte sich ebenfalls auf die \textit{Ephemerides barometricae} beziehen, in denen Ramazzini das "paradoxe" Verhalten des Luftgewichtes sowohl empirisch nachzuweisen wie theoretisch zu verdeutlichen versucht und auch weitere Erklärungen des Phänomens ausfürlich erörtert, welche tatsächlich Leibnizens Ansatz zum Teil ähneln: u.\,a. etwa diejenigen von Joachim d'Alencé\indexN{\textls{Alencé} (d'Alancé, Dalancé, Dalancay, Dalencé), Joachim de 1640\textendash1707} und Claude Comiers\indexN{\textls{Comiers}, Claude, gest. 1693} (vgl. \textsc{Ramazzini}, \textit{Ephemerides barometricae}, 1695,\cite{01516} S.\;31\,f. u.\;38; 1996,\cite{01517} S.\;12\,f. u.\;15).
Trifft diese Vermutung zu, so %
wurde sowohl die Aufzeichnung N.\;\ref{60267} wie vermutlich auch die Aufzeichnung N.\;\ref{53442} erst verfasst, nachdem Leibniz Ramazzinis \textit{Ephemerides barometricae} gelesen hatte.\cite{01516}
Von deren Veröffentlichung dürfte er zum ersten Mal aus Antonio Magliabechis\indexN{\textls{Magliabechi} (Magliabecchi, Magliabecki, Magliabekius, Magliabequius), Antonio 1633\textendash1714} Brief vom 26.\;April 1695 erfahren haben, in dem auch die gelegentliche Sendung eines Exemplars der Abhandlung %
an Leibniz in Aussicht gestellt wurde (\textit{LSB} I,\;11 N.\;294, S.\;424.23--425.16).\cite{01441}
Dem weiteren Briefaustausch beider Gelehrten ist jedoch nicht zu entnehmen, dass diese Sendung jemals stattgefunden habe.
Ramazzini\indexN{\textls{Ramazzini} (Ramazzinus), Bernardino 1633\textendash1714} selbst teilte aber in seinem Brief an Leibniz vom 14.\;August 1696 mit, dem nach Hannover\indexO{Hannover} zurückreisenden Arzt Georg Erich Barnstorf\indexN{\textls{Barnstorf}, Georg Erich 1655\textendash1729} drei Exemplare der Abhandlung\cite{01516} %
für Leibniz anvertraut zu haben (\textit{LSB} III,\;7 N.\;22, S.\;87.10--13).\cite{01442}
Leibnizens Schreiben an Ramazzini\indexN{\textls{Ramazzini} (Ramazzinus), Bernardino 1633\textendash1714} vom 4.\;(14.) September 1696 zeigt allerdings, dass Barnstorf\indexN{\textls{Barnstorf}, Georg Erich 1655\textendash1729} bislang seinen Auftrag nicht erledigt hatte (\textit{LSB} III,\;7 N.\;31, S.\;118).\cite{01446}
Es ist daher anzunehmen, dass Leibniz vor Mitte September 1696 keine Gelegenheit hatte, Ramazzinis\indexN{\textls{Ramazzini} (Ramazzinus), Bernardino 1633\textendash1714} \textit{Ephemerides barometricae}\cite{01516}\cite{01517} zu lesen.
\textit{Mitte September 1696} ist demgemäß als frühester Entstehungszeitpunkt der Aufzeichnung N.\;\ref{60267} zu betrachten.
Dass die Aufzeichnung N.\;\ref{53442} ebensowenig vor diesem Datum verfasst wurde, ist nicht zwingend, aber wahrscheinlich.
\pend%
\pstart%
(3)
Ob Leibniz die erste Ausgabe von Ramazzinis\indexN{\textls{Ramazzini} (Ramazzinus), Bernardino 1633\textendash1714} \textit{Ephemerides barometricae} (Modena 1695)\cite{01516} oder die zweite (Augsburg 1696)\cite{01517} gelesen hat, ist unklar.
Die Augsburger\indexO{Augsburg} Ausgabe wurde erstmals auf der Leipziger\indexO{Leipzig} Ostermesse 1696 ausgestellt und dürfte somit ab dem Frühjahr desselben Jahres auf dem Büchermarkt erhältlich gewesen sein (vgl. \textit{Catalogus universalis}, Ostermesse, Leipzig 1696, S.\;B1\,r\textsuperscript{o}).\cite{02578}
Das oben erwähnte Schreiben an Ramazzini\indexN{\textls{Ramazzini} (Ramazzinus), Bernardino 1633\textendash1714} vom 4.\;(14.) September 1696 zeigt jedoch, dass Leibniz zu jenem Zeitpunkt noch gar kein Exemplar der ersten oder der zweiten Ausgabe besessen hat.
(Am Rande sei bemerkt, dass keine bei der GWLB Hannover\indexO{Hannover} aufbewahrten Exemplare der ersten und der zweiten Ausgabe der \textit{Ephemerides barometricae}\cite{01516}\cite{01517} eindeutige Gebrauchsspuren von Leibnizens Hand aufweisen.)
Ramazzinis\indexN{\textls{Ramazzini} (Ramazzinus), Bernardino 1633\textendash1714} Abhandlung wird zum ersten Mal in dem wohl kurz nach dem 12.\;Oktober 1696 verfassten Brief Leibnizens an Michel Angelo Fardella\indexN{\textls{Fardella}, Michel Angelo 1650\textendash1718} (\textit{LSB} II,\;3 N.\;78, S.\;203.20--204.2):\cite{01443}
\textit{Ramazzinus} \lbrack...\rbrack\ \textit{cum observasset Aeris columnam pluvia imminente leviorem fieri, quod dudum Gerickius\indexN{\textls{Guericke} (Gerickius, Gericke), Otto von 1602\textendash1686} primus virunculi prophetae inventor notaverat, miratus est non immerito, ut ipse in suis eruditis "Ephemeridibus" testatur. Causa nulla alia est, quam quod guttae aquae dum dispersae sunt et in aere sustinentur, cum eo gravitant; ubi vero collectae descendunt, jam aerem non amplius sustinentem exonerant. Proposui eam in rem experimentum in libra peragendum.}
Dieser Bericht zeigt, dass Leibniz um Mitte Oktober 1696 nicht nur Ramazzinis\indexN{\textls{Ramazzini} (Ramazzinus), Bernardino 1633\textendash1714} \textit{Ephemerides barometricae}\cite{01516}\cite{01517} gelesen hatte, sondern die kontraintuitive Veränderung des Luftgewichts bei wechselnder Wetterlage auch in gleicher Weise erklärte wie in der Aufzeichnungen N.\;\ref{53442} und N.\;\ref{60267} (\refpassage{37_04_149-150_149v_causa-1}{37_04_149-150_149v_causa-2}).
Insbesondere die in der Vergangenheitsform ausgedrückte Andeutung, er habe zu diesem Sachverhalt ein Experiment vorgeschlagen, das mithilfe einer Waage durchzuführen sei, rechtfertigt die Annahme, dass zu der Zeit, als Leibniz an Fardella\indexN{\textls{Fardella}, Michel Angelo 1650\textendash1718} schrieb, wenigstens die Aufzeichnung N.\;\ref{53442} bereits vorlag, möglicherweise aber auch die Aufzeichnung N.\;\ref{60267}.
\pend%
\pstart%
(4)
Ähnlich wie im Brief an Fardella\indexN{\textls{Fardella}, Michel Angelo 1650\textendash1718} äußert sich Leibniz in seinem Schreiben an Ramazzini\indexN{\textls{Ramazzini} (Ramazzinus), Bernardino 1633\textendash1714} vom 13.\;(23.) Januar 1697.\cite{01444}
Dort zeigt eine Bemerkung über Rudolf Jakob Camerarius\indexN{\textls{Camerarius}, Rudolf Jakob 1665\textendash1721} zusätzlich, dass Leibniz spätestens zu diesem Zeitpunkt dessen \textit{Ephemerides meteorologicae} (Augsburg 1696)\cite{01445} kannte (\textit{LSB} III,\;7 N.\;67, S.\;264.3--14).\cite{01444}
Ein früherer Hinweis auf Camerarius'\indexN{\textls{Camerarius}, Rudolf Jakob 1665\textendash1721} meteorologisches Werk,\cite{01445} das zugleich die zweite Ausgabe von Ramazzinis\indexN{\textls{Ramazzini} (Ramazzinus), Bernardino 1633\textendash1714} \textit{Ephemerides barometricae}\cite{01517} umfasste (siehe oben), ist im Leibniz-Briefwechsel nicht bekannt.
Camerarius\indexN{\textls{Camerarius}, Rudolf Jakob 1665\textendash1721} wird aber namentlich -- zusammen mit Salomon Reisel\indexN{\textls{Reisel} (Reyselius, Risselius, Reuschius), Salomon 1625\textendash1702} -- auch in einer Ergänzung zur Aufzeichnung N.\;\ref{60267} (\refpassage{LH_37_04_149-150_150v_Camerarius-1}{LH_37_04_149-150_150v_Camerarius-2}) erwähnt.
Dies lässt sich als Zeichen dafür deuten, dass N.\;\ref{60267} etwas später als die Aufzeichnung N.\;\ref{53442} (wo Camerarius unerwähnt bleibt) entstand, etwa um Mitte Januar 1697 wie der zuletzt erwähnte Brief an Ramazzini\indexN{\textls{Ramazzini} (Ramazzinus), Bernardino 1633\textendash1714}.
Eine noch spätere Entstehungszeit -- nach dem Frühjahr 1697 -- ist anzunehmen, wenn Leibniz in der Aufzeichnung N.\;\ref{60267} (\refpassage{LH_37_04_149-150_150v_hllnwtn-1}{LH_37_04_149-150_150v_hllnwtn-2}) tatsächlich auf Edmond Halleys\indexN{\textls{Halley} (Hali, Hallé, Hale, Hales, Hallaeus), Edmond 1656\textendash1742} Aufsatz "The true theory of the tides, extracted from 'hat admired treatise of Mr.\;Isaac Newton\indexN{\textls{Newton} (Neutonus), Isaac 1643\textendash1727} intituled \textit{Philosophiae naturalis principia mathematica}" anspielt, welcher im Heft der \textit{Philosophical transactions} vom 31.\;März 1697 (Bd.\;XIX, Nr.\;226, S.\;445--457)\cite{02593}\cite{00158} veröffentlicht worden war.
Hieraus ergibt sich als möglicher Terminus post quem der Abfassung von N.\;\ref{60267} der \textit{Frühjahr 1697}.
Da diese Schlussfolgerung aber unter Vorbehalt steht, %
ist auch möglich, dass Leibniz den Hauptteil der Aufzeichnung N.\;\ref{60267} %
doch früher verfasst und die besagte Ergänzung über Camerarius\indexN{\textls{Camerarius}, Rudolf Jakob 1665\textendash1721} und Reisel\indexN{\textls{Reisel} (Reyselius, Risselius, Reuschius), Salomon 1625\textendash1702} erst nachträglich in N.\;\ref{60267} hinzugefügt hat.
Hierfür spricht %
der Umstand, dass der gesamte Schlussteil von N.\;\ref{60267} (ab \refpassage{LH_37_04_149-150_150v_Camerarius-1}{LH_37_04_149-150_150v_Camerarius-1}) einen etwas anderen Schreibduktus als der vorausgehende Hauptteil aufweist und sich von diesem letzteren auch inhaltlich abhebt.
Der gesamte Schlussteil konnte somit nachträglich ergänzt worden sein.
\pend%
\pstart%
(5)
Als Leibniz in seinem Brief an Ramazzini\indexN{\textls{Ramazzini} (Ramazzinus), Bernardino 1633\textendash1714} vom 12.\;(22.) April 1699 wieder auf das paradoxe Verhalten des Luftgewichts Bezug nimmt, bezeichnet er die \textit{Ephemerides barometricae}\cite{01516}\cite{01517} als "zeitnah" und nennt sie im gleichen Atemzug mit Camerarius'\indexN{\textls{Camerarius}, Rudolf Jakob 1665\textendash1721} \textit{Ephemerides meteorologicae},\cite{01445} während er seinen eigenen Erklärungsvorschlag als älter darstellt (\textit{LSB} III,\;8 N.\;31, S.\;100.11--14):\cite{01447}
\textit{Cur Aeris pluvii columna levior sit, quod contra accidere debere videbatur, hujus et perelegantis problematis tua cum Experientissimo Camerario aliisque tractatio nupera utique praeclara multa continet; legique cum fructu et voluptate, etsi causam mechanicis nixam demonstrationibus jam dudum exploratam habere mihi sim visus.}
Diese Ausführung %
ist mit der Annahme vereinbar, dass die Aufzeichnung N.\;\ref{60267} ab dem Frühjahr 1697 verfasst wurde;
sie zeigt aber, dass wenigstens der Schlussteil von N.\;\ref{60267} (samt Verweis auf Camerarius\indexN{\textls{Camerarius}, Rudolf Jakob 1665\textendash1721} und Reisel\indexN{\textls{Reisel} (Reyselius, Risselius, Reuschius), Salomon 1625\textendash1702}) noch im April 1699 entstanden sein könnte oder später.
Dies ist ferner mit der Datierung des im Träger von N.\;\ref{60267} (Bl.\;\href{\Scan{LH_37_04_150wz}}{150}) vorliegenden Wasserzeichens vereinbar, dessen Vorkommen im Leibniz-Nachlass nach heutigem Wissensstand bis zum Jahr 1697 belegt ist:
Das Wasserzeichen zeigt zwar, dass wenigstens der Hauptteil von N.\;\ref{60267} -- wenn nicht die gesamte Aufzeichnung -- noch im Jahr 1697 oder etwas später niedergeschrieben worden sein sollte;
einer möglichen Ergänzung des Schlussteils zu einem späteren Zeitpunkt steht es aber nicht im Wege.
Hieraus ergibt sich als wahrscheinlicher Terminus ante quem der Datierung -- wenigstens für den Hauppteil von N.\;\ref{60267} -- die \textit{zweite Hälfte 1697 oder etwas später}.
\pend%
\pstart%
\edlabel{paradoxon_RamaSchel-fhg-1}%
(6)
Wahrscheinlich hatte Leibniz, als er die Aufzeichnung N.\;\ref{60267} samt deren möglicherweise nachträglich ergänztem Schlussteil verfasste, von Günther Christoph Schelhammers\indexN{\textls{Schelhammer} (Schelhammerus), Günther Christoph 1649\textendash1716} \textit{Solutio problematis cur mercurius in tubo Torricelliano, seu barometro, pluvioso tempore descendat cum deberet ascendere} noch keine Kenntnis.\cite{01450}
Denn dieser 1697 in den \textit{Miscellanea curiosa} (decuria III, annus IV, Appendix, S.\;54--62)\cite{01459} veröffentlichte Aufsatz -- Schelhammers\indexN{\textls{Schelhammer} (Schelhammerus), Günther Christoph 1649\textendash1716} erster Beitrag zur Erforschung der wetterbedingten Veränderung der Quecksilberhöhe im Barometer und zugleich der Beginn seines langwierigen Disputs mit Ramazzini -- bleibt in der gesamten, ansonsten an Literaturhinweisen großzügigen Aufzeichnung N.\;\ref{60267} unerwähnt.
Von Schelhammers \textit{Solutio problematis}\cite{01450} dürfte Leibniz erst erfahren haben, nachdem er Ramazzinis\indexN{\textls{Ramazzini} (Ramazzinus), Bernardino 1633\textendash1714} Brief vom 17.\;Juni 1699 empfangen hatte (vgl. \textit{LSB} III,\;8 N.\;51, S.\;153.5--8\cite{01448}).
Beigelegt war nämlich ein Nachdruck von Schelhammers\indexN{\textls{Schelhammer} (Schelhammerus), Günther Christoph 1649\textendash1716} Aufsatz samt Ramazzinis\indexN{\textls{Ramazzini} (Ramazzinus), Bernardino 1633\textendash1714} eigener Erwiderung unter dem gemeinsamen Titel \textit{Epistola G.\:C.\:Schelhammeri de motionibus mercurii in Torricelliana fistula ... Accedit ejusdem argumenti epistola B.\:Ramazzini} (Modena 1698,\cite{01451} Schelhammers Aufsatz auf S.\;1--14;\cite{01450} die Beilage ist im Leibniz-Nachlass unter der Signatur LH~XXXVII~4 Bl.\;172--173 erhalten).
Eine frühere Erwähnung der \textit{Solutio problematis}\cite{01450} im Leibniz-Briefwechsel ist nicht bekannt.
In seiner Antwort an Ramazzini\indexN{\textls{Ramazzini} (Ramazzinus), Bernardino 1633\textendash1714} vom 28.\;Dezember 1699 (7.\;Januar 1700) spielt Leibniz indessen auf Schelhammers zweiten Beitrag über die wetterbedingten Veränderungen des Lufgewichts an (\textit{LSB} III,\;8 N.\;103, S.\;281.13--15).\cite{01449}
Diese erneute Anfechtung von Ramazzinis Ansichten war in einem von Schelhammer selbst herausgegebenen \glqq Sammelband\grqq\ erschienen: \textit{De motu mercurii in tubo Torricelliano epistola. Accessit altera ejusdem argumenti B.\;Ramazzini ... itemque tertia, sive ejusdem Schelhameri ad objectiones Ramazzinianas responsio} (Kiel 1699, S.\;69--150);\cite{01452}
dort war auch seine \textit{Solutio problematis} ein zweites Mal abgedruckt worden (S.\;10--28).\cite{01450}\cite{01452}
(Das bei der GWLB Hannover\indexO{Hannover} unter der Signatur N-A 1193 aufbewahrte Exemplar von \textit{De motu mercurii in tubo Torricelliano}\cite{01452} weist Gebrauchspuren auf, die auf Leibniz zurückgehen könnten.)
Folglich dürfte er wohl zwischen Juli und Dezember 1699 Schelhammers\indexN{\textls{Schelhammer} (Schelhammerus), Günther Christoph 1649\textendash1716} Beiträge zur damaligen Diskussion um das paradoxe Verhalten des Luftgewichts gelesen haben;
dass er sie bereits früher kannte, ist nicht unmöglich, aber mangels Erwähnungen im Briefwechsel unwahrscheinlich.
Daraus lässt sich schließen, dass der Schlussteil von N.\;\ref{60267} -- falls überhaupt nachträglich ergänzt -- spätestenst in der \textit{zweiten Hälfte 1699} abgefasst worden sein düfte.
Dieser Zeitpunkt ist somit auch als spätester Terminus ante quem für die gesamte Aufzeichnung N.\;\ref{60267} zu betrachten.%
\edlabel{paradoxon_RamaSchel-fhg-3}
\pend%
\pstart%
(7)
Die titellose und undatierte Aufzeichnung N.\;\ref{36805}, editorisch \textit{Experimentum mechanicum de paradoxi barometrici causa} benannt, geht indessen gleich zu Beginn (\refpassage{LBr_413_048_048v_Schelh-1}{LBr_413_048_048v_Schelh-2}) auf Schelhammers\indexN{\textls{Schelhammer} (Schelhammerus), Günther Christoph 1649\textendash1716} Beiträge ein.
Die \textit{zweite Hälfte 1699} ist demgemäß als Terminus post quem der Entstehung von N.\;\ref{36805} anzusehen.
Dass Leibniz an der angegebenen Stelle (\refpassage{LBr_413_048_048v_Schelh-1}{LBr_413_048_048v_Schelh-2}) von den \textit{scripta nupera et mutua virorum clarissimorum Camerarii,\indexN{\textls{Camerarius}, Rudolf Jakob 1665\textendash1721} Ramazzini,\indexN{\textls{Ramazzini} (Ramazzinus), Bernardino 1633\textendash1714} Schelhameri\indexN{\textls{Schelhammer} (Schelhammerus), Günther Christoph 1649\textendash1716}} spricht, ist wohl als Zeichen für eine eher zeitnahe Entstehung von N.\;\ref{36805} zu deuten.
Demnach könnte die Aufzeichnung schon vorgelegen haben, als Leibniz am 18.\;März 1700 seine Erklärung der wetterbedingten Veränderung des Luftgewichts an Ramazzini mitteilte (vgl. \textit{LSB} III,\;8 N.\;139, S.\;371.13--372.14);\cite{01454}
um eine solche Mitteilung hatte Ramazzini\indexN{\textls{Ramazzini} (Ramazzinus), Bernardino 1633\textendash1714} selbst in seinem Brief vom 24.\;Februar 1700 gebeten (vgl. \textit{LSB} III,\;8 N.\;130, S.\;336.13--15).\cite{01453}
Leibnizens Brief vom 18.\;März 1700\cite{01454} und die Aufzeichnung N.\;\ref{36805} weisen zweifelsohne eine deutliche inhaltliche Verwandschaft auf, wobei es aber schwierig zu bestimmen ist, welcher Text welchem tatstächlich als Vorlage gedient haben könnte.
\pend%
\pstart%
(8)
Falls die Aufzeichnung N.\;\ref{36805} später entstand, könnte sie im Zusammenhang mit Leibnizens Brief an Friedrich Hoffmann\indexN{\textls{Hoffmann}, Friedrich 1660\textendash1742} vom 19.\;März 1701 stehen, dem eine \textit{explanatio phaenomeni in barometro} beigelegt war (vgl. III,\;8 N.\;225, S.\;577.17\,f.).\cite{01456}
Die Aufzeichnung könnte somit als Vorarbeit zu der nach heutigem Wissensstand verschollenen Beilage für Hoffmann\indexN{\textls{Hoffmann}, Friedrich 1660\textendash1742} gedient haben.
Hieraus ergibt sich der \textit{19.\;März 1701} als spätester Zeitpunkt, zu dem N.\;\ref{36805} angefertigt worden sein dürfte.
Für eine Entstehung der Aufzeichnung in diesem Zusammenhang spricht auch der Umstand, dass der Träger von N.\;\ref{36805} im Leibniz-Nachlass unter der Signatur LBr~413 überliefert ist, d.\,h. zusammen mit weiteren Konzepten von Leibnizens Briefen an Hoffmann.
Für eine frühere Entstehung spricht hingegen, dass Leibniz seine Erklärung bereits in seinem Schreiben an Hoffmann\indexN{\textls{Hoffmann}, Friedrich 1660\textendash1742} vom 16.\;Februar 1701 angekündigt und sogleich in die Zukunft verschoben hatte, weil ihm seine \textit{pristinae schedulae} nicht zu Hande gewesen seien (\textit{LSB} III,\;8 N.\;212, S.\;551.4--8).\cite{01455}
Falls Leibniz damit auch auf die Aufzeichnung N.\;\ref{36805} anspielt, dann ist diese vor dem 16.\;Februar 1701 entstanden:
Als er die Beilage für Hoffmann abfasste, könnte er die schon vorliegende Aufzeichnung zu Rate gezogen und sie dann einfach zusammen mit den übrigen Konzepten der Briefe an Hoffmann\indexN{\textls{Hoffmann}, Friedrich 1660\textendash1742} aufbewahrt haben.
Unter den \textit{pristinae schedulae} könnte er aber lediglich die früheren Aufzeichnungen N.\;\ref{60267} oder N.\;\ref{53442} gemeint haben, oder auch andere nicht ermittelte Texte zum gleichen Thema.
Eine genauere chronologische Einordnung der Aufzeichnung N.\;\ref{36805} ist angesichts dieser Unklarheit nicht möglich.
\pend%
\pstart%
(9)
Die in Leibnizens Brief an Ramazzini vom 18.\;März 1700 übermittelte Erklärung des paradoxen Verhaltens des Luftgewichtes\cite{01454} (siehe oben) hatte beim italienischen Arzt und Naturforscher, der im selben Jahr an die Universität zu Padua\indexO{Padua} berufen worden war, offenbar Beifall gefunden:
Ramazzini\indexN{\textls{Ramazzini} (Ramazzinus), Bernardino 1633\textendash1714} betrachtete Leibnizens Erklärung als richtige Lösung des Paradoxons und adäquate Antwort auf die Kritik, die Schelhammer\indexN{\textls{Schelhammer} (Schelhammerus), Günther Christoph 1649\textendash1716} in seinen polemischen Schriften \textit{Solutio problematis}\cite{01450} und \textit{De motu mercurii}\cite{01452} geübt hatte (siehe oben),
und er plante bereits zu diesem Zeitpunkt eine Erwiderung auf Schelhammer,\indexN{\textls{Schelhammer} (Schelhammerus), Günther Christoph 1649\textendash1716} wie Friedrich August Hackmann\indexN{\textls{Hackmann}, Friedrich August 1670\textendash1742} und Ramazzini selbst in ihren Briefen an Leibniz vom 8.\;Juni 1700 (\textit{LSB} I,\;18 N.\;75, S.\;103.20--21)\cite{01464} bzw. von der ersten Junihälfte 1704 (\textit{LSB} III,\;9 N.\;182 S.\;521.2--5)\cite{01466} festhalten.
Aufgrund einer schweren Erkrankung sollten jedoch insgesamt zehn Jahre vergehen, ehe Ramazzini\indexN{\textls{Ramazzini} (Ramazzinus), Bernardino 1633\textendash1714} wieder zur Feder griff und einen weiteren Beitrag zur Diskussion über die wetterbedingten Schwankungen des Barometeres und zur gelehrten Polemik mit Schelhammer\indexN{\textls{Schelhammer} (Schelhammerus), Günther Christoph 1649\textendash1716} leistete.
Dies erfolgte mit der dritten Auflage seiner meteorologischen Abhandlung:
\textit{Ephemerides barometricae cum disquisitione causae ascensus ac descensus mercurii in Torricelliana fistula juxta diversum aeris statum a Bernardino Ramazzini ... recusa cum tota controversia quam idem habuit cum C.\:G.\;Schelhamero ... Hisce omnibus jam editis accedit nova epistola ejusdem Ramazzini ad L.\;Schroekium ... cum solutione problematis ex invento V.\:G.\;Leibnitii} (Padua 1710).\cite{01462}
Wie der ausführliche Titel der neuen Ausgabe andeutet, umfasste der dort veröffentlichte Brief Ramazzinis\indexN{\textls{Ramazzini} (Ramazzinus), Bernardino 1633\textendash1714} (vom 1.\;Juni 1710) an Lucas Schröck auch einen Abdruck von Leibnizens Brief an Ramazzini vom 18.\;März 1700 (S.\;196\,f.).\cite{01460}\cite{01454}
Anschließend berichtete Ramazzini\indexN{\textls{Ramazzini} (Ramazzinus), Bernardino 1633\textendash1714} im gleichen Brief an Schröck,\indexN{\textls{Schröck} (Schroekius), Lucas 1646\textendash1730} mit Unterstützung seines Kollegen Giovanni Graziani\indexN{\textls{Graziani} (Gratianus), Giovanni, gest. 1744} ein Experiment durchgeführt zu haben, welches das von Leibniz vorgeschlagene nachgeahmt und empirisch bestätigt habe (ebd. S.\;197--201).\cite{01462}\cite{01460}
Ramazzinis\indexN{\textls{Ramazzini} (Ramazzinus), Bernardino 1633\textendash1714} Brief an Schröck\indexN{\textls{Schröck} (Schroekius), Lucas 1646\textendash1730} wurde ferner -- unter dem Titel \glqq Epistola de causa adscensus et descensus mercurii in fistula Torricelliana\grqq\ -- in den \textit{Academiae naturae curiosorum ephemerides}, centuriae I--II, 1712 abgedruckt (Anhang, Nr.\;3, S.\;92--98;\cite{01457} Leibnizens Brief findet sich dort auf S.\;95,\cite{01454} Ramazzinis Bericht über das Experiment auf S.\;95--97).\cite{01460}
Im selben Band der \textit{Academiae naturae curiosorum ephemerides} veröffentlichte auch Schelhammer\indexN{\textls{Schelhammer} (Schelhammerus), Günther Christoph 1649\textendash1716} seine kritischen Bemerkungen zu Leibnizens Erklärung und Ramazzinis\indexN{\textls{Ramazzini} (Ramazzinus), Bernardino 1633\textendash1714} Experiment (\glqq Schelhameri notata ad hanc epistolam\grqq, Anhang, Nr.\;4, S.\;98--106).\cite{01461}\cite{01457}
Dort stellt Schelhammer\indexN{\textls{Schelhammer} (Schelhammerus), Günther Christoph 1649\textendash1716} insbesondere in Frage, ob Ramazzinis\indexN{\textls{Ramazzini} (Ramazzinus), Bernardino 1633\textendash1714} empirischer Versuch Leibnizens Hypothese wirklich bestätig habe (ebd. S.\;98\,f.).\cite{01461}\cite{01457}%
\edlabel{paradoxon_RamaSchel-fhg-2}%
\pend%
\pstart%
(10)
Leibniz hatte diese kritischen Bemerkungen offenbar gelesen, als er am 17.\;Juni 1712 in einer knappen Mitteilung an Schelhammer\indexN{\textls{Schelhammer} (Schelhammerus), Günther Christoph 1649\textendash1716} darauf Bezug nahm (\textit{LOD}~II.2, S.\;70).\cite{01463}
Auf die Einwände des Kieler Arztes und Hochschullehrers geht Leibniz ausführlicher in der Erläuterung N.\;\ref{54958} \textit{De aetiologia barometri} ein, die er in der folgenden Ausgabe der \textit{Academiae naturae curiosorum ephemerides} von 1715 veröffentlichen ließ (Appendix, Nr.\;3, S.\;49--51).\cite{01457}\cite{01458}
Es ist anzunehmen, dass er diese kurze Stellungnahme N.\;\ref{54958}, zu der nach heutigem Kenntnisstand keine unmittelbaren Vorarbeiten überliefert sind, erst nach der Mitteilung an Schelhammer vom 17.\;Juni 1712 verfasste.
Die \textit{zweite Hälfte Juni 1712} gilt somit als Terminus post quem der Datierung von N.\;\ref{54958}.
Die (nicht aufgefundene) handschriftliche Vorlage für den Druck entstand wohl spätestens im Laufe des Jahres 1714 oder in den ersten Monaten des folgenden Jahres, jedenfalls rechtzeritig für die Veröffentlichung.
Wann genau im Jahre 1715 die betreffende Ausgabe der Zeitschrift der Leopoldina-Akademie ausgeliefert wurde, ist nicht bekannt.
Es ist aber wahrscheinlich, dass Leibniz den Text der Erläuterung N.\;\ref{54958} schon bald nach seiner Mitteilung an Schelhammer\indexN{\textls{Schelhammer} (Schelhammerus), Günther Christoph 1649\textendash1716} vom 17.\;Juni 1712 verfasst hatte.\cite{01463}
\pend%
\pstart%
Insgesamt ist der Textkomplex N.\;\ref{paradoxon_intro} \textit{Paradoxon de aeris gravitate} auf den Zeitraum von Mitte September 1696 bis zum Jahr 1715 zu datieren.
Die der Reihe nach älteste Aufzeichnung N.\;\ref{53442} könnte auch früher (etwa ab 1689) entstanden sein und wurde wahrscheinlich nicht nach Mitte Oktober 1696 angefertigt.
Die sicherlich nicht vor Mitte September 1696 und vermutlich nicht viel später als N.\;\ref{53442} entstandene Aufzeichnung N.\;\ref{60267} könnte in der zweiten Hälfte 1697 oder etwas später verfasst worden sein, wohl aber nicht nach der zweiten Hälfte 1699, als ihr Schlussteil spätestens hinzugefügt wurde (falls Leibniz ihn überhaupt nachträglich ergänzt hat).
Die Aufzeichnung N.\;\ref{36805} entstand nicht vor der zweiten Hälfte 1699 und wahrscheinlich nicht nach dem 19.\;März 1701.
Die im Laufe des Jahres 1715 veröffentlichte Stellungnahme N.\;\ref{54958} \textit{De aetiologia barometri} wurde wahrscheinlich nicht vor der zweiten Hälfte Juni 1712 abgefasst.
\pend%
\selectlanguage{latin}%
\frenchspacing%
\normalsize%
